保持风扇设置不变，将加热电压阶跃到 3 V，并记录约 240 s。图 2 显示，T2 因靠近加热元件而升温最快、幅值最大；T3 位于出口，响应更慢且幅值较小。T1 也有明显变化，说明入口空气温度并非常数。

从加热电压到 T2 的主要动态可用 PT1 环节近似；从输入到 T3 还包含由 T2 到出口的热传递和输运延迟。实际控制器设计采用 b.m 中给出的 PT1 模型。由于 T1 在测量中变化较大，而且现有文件无法完全分离所有测量序列，因此不再从初值、终值强行重新估计两个对象，避免得到不可靠参数。

Simulink 模型按照“加热电压—T2—T3”的信号关系搭建。闭环中先计算设定值与输出的误差，然后依次经过控制器和 PT1 对象。末端的 25 °C 常数用于把增量模型转换为绝对温度，反馈为负反馈。

严格比较要求把实测与仿真的 T2、T3 放在同一图中，并使用相同初值和 3 V 阶跃。当前文件没有导出的开环仿真曲线，因此只能定性比较：模型能描述主要的单调升温过程，但实测曲线还受到 T1、空气流量、传感器滞后和可能的纯延迟影响。
4.3 长时间测量（c）
任务要求读取 Sprung_3V_Langzeit.mat 并与本组 240 s 测量比较，但该文件没有单独提供。因此本节仅说明规范的比较方法，不给出未经验证的数值。应重点比较 240 s 时是否真正达到稳态、两次实验的初始温度和空气流量是否一致，以及 T1 是否具有相同变化趋势。若长期终值更高，则用 240 s 数据计算的静态增益会被低估。
4.4 输入电压—T2 静态特性（d）
应从 KL_U_T.mat 中提取 0 至 5 V、步长 0.5 V 的各稳态 T2 值，绘制 T2,stat=f(U)。当前没有单独的 KL_U_T.mat，因此没有可验证的完整数值表，本报告不填入未经验证的数据。该特性用于识别近似线性区、低电压死区和高电压饱和，并限定线性控制器适用的工作点范围。
4.5 5 V 测量与最短升温时间（e）
任务要求利用 sprung_5_V.mat 判断最大输入下达到指定温度的最短时间。该数据文件未单独提供，所以无法给出具体秒数。可由物理可实现性判断：6 s 的调节时间对于该热系统明显过快；100 s 与现有响应和仿真相符；300 s 虽然容易实现，但过于保守。因此后续选择 100 s 量级的期望动态。
4.6 期望传递函数和控制器（f）
最大允许超调量为 5%，取阻尼比 D=0.69。脚本中选择 ωn=0.055 rad/s，对应约 100 s 量级的调节时间。期望闭环模型为：
G_M(s) = 0.00308 / (s² + 0.0766 s + 0.00308)	(7)
把该模型和 PT1 对象代入代数控制器公式并约分，得到：
K(s) = (0.004527 s + 0.0001886) / (s² + 0.0766 s)	(8)
分母含因子 s，因此控制器具有积分环节。另一个极点和零点用于整形，使闭环近似满足所要求的 PT2 动态。
4.7 闭环 Simulink 仿真（g）
仿真以 25 °C 为温度偏置，目标温度约为 28 °C。输出平稳上升，峰值约 28.1 °C，随后收敛。相对于 3 K 的设定值变化，超调量约为 3% 至 5%；约 100 至 150 s 后输出接近最终值。

4.8 dSpace 实物实验（h）
控制器应在 Simulink 中接入 dSpace 模型，设置合理的 0 至 5 V 限幅，并同步记录设定值、T1、T2、T3 和控制电压。当前测量文件中没有一个能明确识别为完整闭环试验的独立数据集，因此不能做定量的仿真—实物对比。本节据实验流程说明实施步骤，不填入缺乏原始数据支持的实测数值。
T1 是解释偏差的关键变量。T1 上升时，进入管道的空气已经变暖，T2 和 T3 可能比恒定工作点模型升得更快或更高；T1 下降时则相反。空气流量波动还会改变换热强度和输运时间。
4.9 无积分环节控制器（i）
对于稳定 PT1 对象，去除积分环节后通常不能获得相同稳态品质。P 或 PD 控制器必须依靠非零误差产生恒定控制作用，因此对阶跃设定值一般存在稳态误差。积分环节累积误差，使稳定闭环对恒定设定值的误差趋于零，也能更有效地抑制恒定扰动。无积分控制器可能更不易出现积分饱和，但不能保证准确到达设定值。
